\section{课程概览与深度学习的进步}

MIT 6.S191 是 MIT 官方的深度学习入门课程，由 Alexander Amini 和 Ava Amini 主讲。这是一门为期一周的密集训练营（boot camp），涵盖深度学习的核心理论与实践。2025 年是该课程开设的第 8 年，已累计教授超过 1300 万名全球学生。

\subsection{深度学习的十年进化}

Amini 教授以一个生动的对比开场：仅仅十年前（2015 年），最先进的深度学习人脸生成系统只能产出模糊的像素块（Goodfellow et al.）；到 2018 年（Karras, Laine, Aila），生成的人脸已具备照片级真实感；2020 年，MIT 6.S191 课程组更是制作了一段以"奥巴马"形象欢迎学员的 deepfake 视频，该视频在 YouTube 上获得了超过 110 万次观看。

\begin{figure}[H]
\centering
\includegraphics[width=0.85\textwidth]{slides-images/slide-002.jpg}
\caption{深度学习人脸生成技术的十年进步：从 2015 年的模糊像素到 2020 年的逼真视频\protect\footnotemark}
\end{figure}
\footnotetext{Slides 第 2 页。}

\begin{knowledgebox}{2020 年 deepfake 的制作成本}
2020 年制作那段 2 分钟的 deepfake 视频需要：2 小时的专业音频录制、50 小时的高清视频数据、预定义的静态脚本、以及超过 15000 美元的计算资源。而到 2025 年，Amini 现场演示了实时语音克隆：仅录制数分钟讲课音频，即可立刻生成他的声音克隆体，并与之进行无脚本的动态对话。这充分说明了生成式 AI 在短短几年内的巨大飞跃。
\end{knowledgebox}

\begin{figure}[H]
\centering
\includegraphics[width=0.85\textwidth]{slides-images/slide-006.jpg}
\caption{2020 年 vs 2025 年：生成式 AI 能力的对比\protect\footnotemark}
\end{figure}
\footnotetext{Slides 第 6 页。}

\subsection{什么是深度学习？}

在理解深度学习之前，需要先厘清三个层次的概念：

\begin{importantbox}{AI $\supset$ ML $\supset$ DL 的层级关系}
\begin{itemize}
\item \textbf{Artificial Intelligence (AI)}：使计算机模拟人类行为的任何技术
\item \textbf{Machine Learning (ML)}：不显式编程，而是让机器从数据中学习模式的方法
\item \textbf{Deep Learning (DL)}：使用深层神经网络从原始数据中提取模式的 ML 子集
\end{itemize}
核心理念：\textbf{教会计算机直接从原始数据中学习如何完成任务}（learn a task directly from raw data）。
\end{importantbox}

\begin{figure}[H]
\centering
\includegraphics[width=0.85\textwidth]{slides-images/slide-007.jpg}
\caption{AI、Machine Learning 与 Deep Learning 的嵌套关系\protect\footnotemark}
\end{figure}
\footnotetext{Slides 第 7 页。}

\textbf{Intelligence}（智能）的本质是处理信息以指导未来决策的能力。Artificial Intelligence 就是构建人工算法来完成同样的过程——利用数据来驱动决策。Machine Learning 不直接告诉计算机如何处理数据，而是让它从数据中自主发现模式。Deep Learning 则进一步使用深层神经网络来完成这一过程。

\subsection{课程结构}

MIT 6.S191 为期一周，包含以下模块：

\begin{table}[H]
\centering
\begin{tabular}{lll}
\toprule
\textbf{日期} & \textbf{讲座主题} & \textbf{实验} \\
\midrule
Day 1 & Introduction to Deep Learning & Lab 1: 音乐生成 \\
Day 2 & Deep Sequence Modeling & Lab 2: 计算机视觉 \\
Day 3 & Deep Computer Vision & Lab 3: 大语言模型 \\
Day 4 & Deep Generative Modeling & -- \\
Day 5 & Deep Reinforcement Learning & Final Project \\
\bottomrule
\end{tabular}
\caption{课程日程概览}
\end{table}

今年的亮点包括：同时支持 TensorFlow 和 PyTorch 两个框架的实验，以及全新的 LLM 实验（微调 20 亿参数的 Gemma 模型并构建 AI 评判器）。

\subsection{本章小结}

深度学习在过去十年取得了惊人进步：从模糊的像素生成到实时语音克隆和动态对话。MIT 6.S191 课程旨在通过一周密集训练，让学生掌握驱动这些进步的基础技术。整个课程围绕一个核心理念：教会计算机直接从原始数据中学习任务。

